% Bagian: computability
% Bab: computability-theory
% Subbagian: computable-sets

\documentclass[../../../include/open-logic-section]{subfiles}

\begin{document}

\olfileid[id]{cmp}{thy}{cps}
\olsection{Himpunan Dapat Dihitung}

Kita dapat memperluas gagasan komputabilitas dari fungsi dapat dihitung ke
himpunan dapat dihitung:

\begin{defn}
  Misalkan $S$ suatu himpunan bilangan asli. Maka $S$ \emph{dapat dihitung}
  jika dan hanya jika fungsi karakteristiknya~$\Char{S}$ dapat dihitung. Dengan
  kata lain, $S$ dapat dihitung jika dan hanya jika fungsi
\[
\Char{S}(x) =
\begin{cases}
1 & \text{jika $x \in S$} \\
0 & \text{selain itu}
\end{cases}
\]
dapat dihitung. Demikian pula, suatu relasi $R(x_0, \dots, x_{k-1})$ dapat
dihitung jika dan hanya jika fungsi karakteristiknya dapat dihitung.

Himpunan dan relasi dapat dihitung juga disebut \emph{dapat diputuskan}.
\end{defn}

\begin{explain}
Perhatikan bahwa sekarang kita mempunyai sejumlah gagasan komputabilitas:
untuk fungsi parsial, fungsi, dan himpunan. Jangan mencampuradukkannya!{}
\iftag{TMs}{Mesin Turing yang menghitung suatu fungsi parsial mengembalikan
  keluaran fungsi itu untuk nilai-nilai masukan tempat fungsi tersebut
  terdefinisi; mesin Turing yang menghitung suatu himpunan mengembalikan $1$
  atau~$0$ setelah memutuskan apakah nilai masukan berada dalam himpunan itu
  atau tidak.}{}
\end{explain}

\end{document}

